\documentclass[12pt]{article}

\usepackage[margin=1in]{geometry}
\usepackage{xcolor}
\usepackage{enumitem}
\usepackage{booktabs}
\usepackage{microtype}
\usepackage{parskip}
\usepackage{titlesec}
\usepackage{mdframed}
\usepackage{amsmath}
\usepackage{natbib}
\usepackage[hidelinks]{hyperref}

% ---- Color scheme -----------------------------------------------------------
\definecolor{respblue}{rgb}{0.7,0,0}   % dark red for responses
\definecolor{quotegray}{rgb}{0.35,0.35,0.35}

% ---- Commands ---------------------------------------------------------------
\newcommand{\resp}[1]{\textcolor{respblue}{#1}}

\newenvironment{reviewercomment}{%
  \begin{mdframed}[backgroundcolor=gray!10,linecolor=gray!40,
                   leftmargin=0pt,rightmargin=0pt,
                   innerleftmargin=8pt,innerrightmargin=8pt,
                   innertopmargin=6pt,innerbottommargin=6pt]
  \itshape\color{quotegray}%
}{%
  \end{mdframed}%
}

\newenvironment{response}{%
  \vspace{4pt}%
  \color{respblue}%
  \noindent\textbf{Response:}\space
}{%
  \vspace{10pt}%
}

\setlength{\parskip}{6pt}
\setlength{\parindent}{0pt}

% ---- Title ------------------------------------------------------------------
\title{\textbf{Response to Reviewer 2}\\[4pt]
       \large Automated Segmentation and Classification of Coral Reef Imagery
       Using Deep Learning}
\author{}
\date{September 2026}

\begin{document}
\maketitle
\thispagestyle{empty}

We thank the reviewer for a thorough and constructive read. Below we address
each comment in turn. Our responses appear in \resp{red}; all section, table,
and figure numbers refer to the revised manuscript (\texttt{submission/main.tex}).

\bigskip
\noindent\rule{\linewidth}{0.4pt}
\bigskip

% =============================================================================
\section*{Comment 1 --- Novelty framing}

\begin{reviewercomment}
Novelty framing overstated (``first'' semantic segmentation approach for
colony-scale bleaching assessment) given CoralSCOP and related literature.
Clarify that the novelty is in the combination of SAM2 + bleaching status +
morpho-taxonomic grouping + ASV top-down imagery, not a broad first-in-field
claim.
\end{reviewercomment}

\begin{response}
\textit{The specific novelty claim has been narrowed}  The Related Work paragraph (Section 2) was
rewritten to situate the contribution relative to CoralSCOP
\citet{zheng2024} and other recent work.  The original broad ``first
semantic segmentation pipeline for coral reef imagery'' framing has been removed.
The claim is now scoped specifically to the combination of features not present
in any single prior system: SAM2-based mask proposal tuned for multi-scale coral
detection, an ensemble CNN classifier for morpho-taxonomic class \emph{and}
bleaching status simultaneously, and application to top-down ASV imagery of a
bleaching event.  The Related Work closing sentence now reads: ``Here we present
(1) a novel pipeline combining SAM2 mask proposals with an ensemble CNN
classifier for both morpho-taxonomic class and bleaching status, and (2) the
first application of semantic segmentation to automate colony-scale bleaching
classification from top-down ASV imagery'' --- narrowing the novelty claim to
the specific combination not addressed by prior work.
\end{response}

% =============================================================================
\section*{Comment 2 --- Data / annotation methodology}

\begin{reviewercomment}
Data section: annotations were SAM2-assisted, which may bias evaluation (since
the pipeline itself relies on SAM2).  Report inter-annotator agreement, QC
procedures, correction rules, and any independent expert audit.  Clarify the
``non-coral'' class definition (negative masks weren't explicitly annotated;
they become a major training component via synthetic generation).
\end{reviewercomment}

\begin{response}
\textit{Partially addressed.}  The revised manuscript adds an explicit
description of the synthetic non-coral generation procedure (Section 3.2.3):
a candidate SAM2 mask is assigned a non-coral label whenever its pixel overlap
with any annotated coral polygon is at or below a fixed tolerance
(\texttt{TOLERANCE = 0.2}), i.e.\ at most 20\% of the candidate mask's pixels
may be genuine coral.  A Limitations paragraph (Section 5) now names this as a
source of label noise and frames the 0.2 threshold explicitly as a
bias--variance tradeoff.

\textbf{Remaining gaps:}
\begin{itemize}[topsep=2pt,itemsep=2pt]
  \item Inter-annotator agreement statistics and a description of the expert
    correction workflow have not been added.  We acknowledge this is a
    meaningful gap: SAM2-assisted annotation could in principle inflate
    apparent segmentation quality if the model's own mask boundaries inform
    the ground truth.  A short methods paragraph on QC/correction procedures
    and, where available, an inter-annotator agreement metric (e.g.\
    Cohen's $\kappa$ or mask-level IoU between annotators) should be added
    before final submission.
  \item An independent expert audit of a randomly sampled subset of
    annotations has not been conducted or reported.
\end{itemize}
\end{response}

% =============================================================================
\section*{Comment 3 --- Experimental design / baselines / ablations}

\begin{reviewercomment}
Experimental design under-validated: single reef system, single ASV platform,
one camera geometry, held-out test set from the same domain.  Add baseline
comparisons (SAM2 alone, CoralSCOP, U-Net/DeepLab, Mask R-CNN, non-ensemble
ResNet) and ablations (dual SAM2 config, synthetic non-coral generation,
oversampling, augmentation, ensemble weighting, mask consolidation).
\end{reviewercomment}

\begin{response}
\textit{Partially addressed.}  Two baselines have been added to the revised
manuscript:

\begin{itemize}[topsep=2pt,itemsep=2pt]
  \item \textbf{CoralSCOP} \citep{zheng2024}: run on the same 552
    images and scored with the identical pixel-level IoU/Dice definitions
    (Section 4.1, Table 7).  NeuralReefer outperforms CoralSCOP on every
    metric in both splits; the comparison also includes an IoU exceedance
    curve (Figure 14).
  \item \textbf{Single CNN (``Single CNN'' in Table 7)}: the 5-model ensemble
    replaced with submodel 4 alone (otherwise identical pipeline), providing a
    non-ensemble ResNet baseline.  This isolates the ensemble's contribution
    from the SAM2 mask proposal and consolidation stages.
\end{itemize}

\textbf{Remaining gaps:}  SAM2-alone (no CNN classifier), U-Net/DeepLab,
Mask R-CNN, and the full suite of ablations (dual SAM2 config, synthetic
non-coral generation strategy, augmentation, oversampling, mask consolidation
hyperparameters) have not been added.  We acknowledge these would substantially
strengthen the experimental validation, noting that each ablation requires a full retraining run.
\end{response}

% =============================================================================
\section*{Comment 4 --- Train/validation/test terminology and leakage}

\begin{reviewercomment}
Train/validation/held-out/test terminology is confusing; possible
over-optimistic reporting.  Abstract's 90.9\%/95.8\%/94.3\% appear to be from
Table 4's augmented validation set (1,762 masks), while ecological evaluation
uses the 111 held-out images.  Explicitly state that no masks, images,
augmentations, or hyperparameters from the final 111-image test set were used
anywhere upstream (SAM2 tuning, CNN training, ensemble weighting, threshold
selection).
\end{reviewercomment}

\begin{response}
\textbf{Fully addressed.}  This comment received the most targeted set of
changes in the revision:

\begin{enumerate}[topsep=2pt,itemsep=3pt]
  \item \textbf{Abstract rewritten.}  The abstract now leads with
    pipeline-level live coral cover IoU (0.674) and Dice (0.781) on the
    held-out test set as the primary performance claim, followed by a
    comparison against CoralSCOP (IoU 0.487/Dice 0.614).  The mask-level
    classification figures (90.3\% accuracy, 94.6\% recall, 93.0\% precision)
    are reported as a clearly labeled secondary result---``at the mask level,
    \emph{once a candidate mask has been proposed}, the classification ensemble
    attains\ldots''---so the reader cannot conflate the two quantities.

  \item \textbf{Explicit no-leakage statement.}  The data-partitioning
    paragraph in Section 3.1 now includes: ``The 111 held-out test images were
    withheld entirely from all upstream steps: SAM2 configuration search, CNN
    training, ensemble weight estimation, and threshold selection were each
    conducted on the 441 training images and a separate mask-level validation
    set drawn from the same 441 images only.''  A stacked-generalization
    justification (Wolpert 1992) is cited for the disjoint
    submodel/ensemble/test-set partitioning.

  \item \textbf{Table 4 scope clarified.}  The Table 4 note now reads: ``All
    metrics in this table are mask-level classification statistics; pixel-level
    segmentation quality is reported separately in Table 6.''  The body of
    Section 4.1 makes the causal point explicit: a classification ensemble
    that is accurate on proposed masks (Table 4) cannot recover coral that
    SAM2 never proposed in the first place, so pipeline-level IoU/Dice is
    treated as the primary end-to-end quality measure.
\end{enumerate}
\end{response}

% =============================================================================
\section*{Comment 5 --- Method description ambiguities}

\begin{reviewercomment}
Method description has logical ambiguities:
\begin{itemize}
  \item Section 3.2.1's mask consolidation (uses classification confidence)
    creates an apparent ordering conflict with Figure 4 (implies classify
    $\to$ reject non-coral $\to$ consolidate) vs.\ Sections 3.2.1/3.3 (implies
    consolidate $\to$ classify).  Needs an unambiguous rewrite.
  \item Ensemble aggregation: non-coral class weight not reported.
  \item Optimization details incomplete.
  \item Convexity claim needs justification.
\end{itemize}
\end{reviewercomment}

\begin{response}
\textbf{Fully addressed across multiple passes.}

\begin{enumerate}[topsep=2pt,itemsep=3pt]
  \item \textbf{Pipeline ordering.}  Figure 4 was
    already correct.  The prose in Sections 3.2.1 and 3.3 (``Coral Cover
    Estimation'') had incorrectly described the reverse order and has been
    rewritten to match the code and figure, with an explicit note that the
    subsections' narrative grouping (segmentation $\to$ classification) differs
    from the execution order.

  \item \textbf{Non-coral weight.}  $\nu_{\mathrm{noncoral}} = 0.25$ is now
    reported and justified in Section 3.2.4: the value matches an \emph{a
    priori} prior of roughly 25\% coral cover across the survey, boosting
    recall for the minority coral classes.  A precise argument (citing Lipton,
    Elkan \& Narayanaswamy 2014 on $F_\beta$-optimal Bayes classifiers) is
    provided showing that $\nu_{\mathrm{noncoral}} = 1/4$ is consistent with
    $F_2$-optimal classification ($\beta^2 = 4$ recall-weighting).

  \item \textbf{Optimization details.}  A full Expectation-Maximization /
    Minorize-Maximize derivation for the ensemble weight updates has been added
    as Appendix C, including the complete E-step, M-step, closed-form
    multiplicative update for $\mathbf{W}$, and renormalization step.  A
    summary E-step/M-step description with an MM-surrogate interpretation
    remains in the main text (Section 3.2.5); the full six-equation derivation
    is deferred to the appendix to preserve readability.

  \item \textbf{Convexity claim.}  The original ``novel convex ensemble
    optimization scheme'' framing has been removed.  Section 3.2.5 now
    explicitly states that the joint $(\boldsymbol\alpha, \mathbf{W})$
    optimization is non-convex; only the individual conditional updates
    (holding one set of parameters fixed) are convex.  The contribution is
    reframed around the interpretability of the learned responsibility
    decomposition.
\end{enumerate}
\end{response}

% =============================================================================
\section*{Comment 6 --- Conservative interpretation of results; IoU/Dice/CIs}

\begin{reviewercomment}
Results should be interpreted more conservatively: pixel accuracy $>$90\% can
be inflated by dominant background/easy non-coral pixels.  Add IoU, Dice, F1,
per-class P/R, macro-averaged metrics, confidence intervals.  Treat
LCC/BCC underestimation as a substantive limitation, not a minor residual bias.
\end{reviewercomment}

\begin{response}
\textbf{Substantially addressed.}

\begin{enumerate}[topsep=2pt,itemsep=3pt]
  \item \textbf{New Section 4.1 (``Segmentation Quality: IoU and Dice'').}
    This section now serves as the primary results section for end-to-end
    pipeline quality, explicitly noting that pixel accuracy can be inflated by
    the dominant non-coral background and that IoU/Dice are the appropriate
    summary statistics.

  \item \textbf{Table 6 (per-taxonomy pixel IoU/Dice/precision/recall).}
    Reports pixel-level IoU, Dice, precision, and recall for all 10 taxonomic
    classes (LCC, bleached, and 8 genera), macro-averaged over the 111 test
    images.  Standard errors ($\mathrm{SD}/\!\sqrt{n}$, where $n = 111$
    images) are reported for every cell in a separate column.

  \item \textbf{Table 7 + Figure 14 (CoralSCOP baseline and IoU exceedance
    curve).}  Both NeuralReefer and the Single CNN ablation are compared against
    CoralSCOP on live coral cover IoU/Dice/precision/recall, in-sample and
    out-of-sample, with standard errors throughout (Table 7).  An IoU
    exceedance curve (Figure 14) shows the full distribution of per-image IoU
    for both methods: 81.1\% of NeuralReefer's test predictions reach
    IoU $\geq 0.5$ vs.\ 50.5\% for CoralSCOP.

  \item \textbf{LCC/BCC underestimation reframed.}  The bias decomposition
    discussion (Table 5, Section 4) now treats the observed underestimation as
    a substantive limitation driven by the segmentation bottleneck.  Section 4.1 makes the mechanism explicit: SAM2's failure
    to propose a mask over a colony propagates into every downstream metric
    regardless of classification accuracy.

  \item \textbf{Standard errors vs.\ bootstrap CIs.}  Standard errors
    ($\mathrm{SD}/\!\sqrt{n}$) are reported in Tables 6, 7, and A.8.  The
    reviewer's original request was for confidence intervals, which can be constructed from the SEs.  If the reviewer prefers
    bootstrap CIs specifically, we are happy to add them.
\end{enumerate}
\end{response}

% =============================================================================
\section*{Comment 7 --- Figure quality}

\begin{reviewercomment}
Figure quality: Figures 4, 5 too dense/small text.  Figure 13 panels hard to
read.  Figure 15 negative depth values need clarification.  Figure 1 needs
scale/depth/quality context.  Figure 9 needs clearer legends.  Table 1 should
clarify whether ``average coral cover'' includes both healthy+bleached;
Table 4 should define whether P/R are coral-vs-noncoral, macro-averaged, or
mask-level aggregated.
\end{reviewercomment}

\begin{response}
\textbf{Substantially addressed.}

\begin{enumerate}[topsep=2pt,itemsep=3pt]
  \item \textbf{Figures 4 and 5.}  Both converted to full-page landscape
    (\texttt{sidewaysfigure}) on their own dedicated pages, substantially
    increasing text legibility.  All figures in the paper now share a unified
    theme with a serif Computer Modern font, 40--60pt base
    text, and a consistent panel background.

  \item \textbf{Figure 13.}  The 12-panel density-plot figure was removed and
    replaced by Table 6 (per-class IoU/Dice), which conveys the same
    per-taxonomy information more precisely and compactly.

  \item \textbf{Figure 15 / depth sign convention.}  A new Appendix B
    (``Segmentation Accuracy vs.\ Water Depth'') explains the sign convention
    explicitly: the x-axis is depth \emph{below the water surface} (0 =
    surface, negative = deeper), not distance to the reef substrate.  The
    figure has been converted to a two-panel design (Full Range + Zoomed to
    0--7.5m depth and 75--100\% accuracy), and the caption and appendix prose
    both state the convention clearly.

  \item \textbf{Figure 1.}  The caption now states the pixel dimensions of
    each panel ($1024 \times 1024$ px) and the survey's camera depth-below-
    surface range (1.6--18.7 m).  Physical scale bars and ground sample
    distance cannot be reported: Yellowfin's metadata records depth below the
    water surface, not altitude above the reef substrate, so no
    camera-to-substrate distance is available to compute a GSD.  This
    limitation is noted explicitly in the caption.  Additionally, Figure 1 was
    rebuilt as a true scatter-of-images plot to show spatial and temporal context: panels now sit
    at their literal acquisition time (x-axis) and logged camera depth (y-axis)
    rather than a fixed $3\times3$ grid, giving the reader spatial and temporal
    context for the nine examples.

  \item \textbf{Figure 9 legends.}
    shared color legend was added beneath each row.

  \item \textbf{Table 4 metric definitions.}  The table note now reads: ``All
    metrics are mask-level classification statistics: a mask is a true positive
    if it is predicted as coral (any class), regardless of genus.  Precision
    and recall are therefore coral-vs-noncoral, not macro-averaged across
    genera.''  Table 4 is explicitly scoped to mask-level classification; pixel
    segmentation quality is in Table 6.

  \item \textbf{Table 1 coral cover definition.}  LCC is defined in the paper to include both healthy and bleached coral (Section 3.1).
\end{enumerate}
\end{response}

% =============================================================================
\section*{Comment 8 --- Related Work citations}

\begin{reviewercomment}
Related Work lacks recent underwater/UAV monitoring literature: cites specific
papers to add (Holothurians monitoring UAV 2025; Mussel-YOLO EI 2025;
RBL-YOLO MPB 2024; multi-label coral condition monitoring MPB 2026;
multi-scale coral condition/Drupella MPB 2026; Cost-Effective AUAV monitoring
MER 2026).
\end{reviewercomment}

\begin{response}
\textit{Not yet addressed.}  The six papers named by the reviewer have not been
added to the Related Work section.
\end{response}

% =============================================================================
\section*{Comment 9 --- Copyediting}

\begin{reviewercomment}
Numerous typos (``heatstress,'' ``stake holders,'' ``over fishing,'' ``land
based,'' ``long term,'' ``limited shallower reefs,'' ``photodata,'' and others)
and an overly informal acknowledgement.  Also: ``Beijorn'' should be
``Beijbom,'' ``Marcohorn'' should be ``Macrohon,'' SAM2 is cited with the
original SAM reference in places, and CoralMask/USIS10K/CoralSCOP references
appear mismatched or duplicated.
\end{reviewercomment}

\begin{response}
\textbf{Substantially addressed.}  All individually named typos have been
corrected in the revised manuscript:

\begin{itemize}[topsep=2pt,itemsep=1pt]
  \item ``heatstress'' $\to$ ``heat stress''
  \item ``stake holders'' $\to$ ``stakeholders''
  \item ``over fishing'' $\to$ ``overfishing''
  \item ``land based'' $\to$ ``land-based''
  \item ``long term'' $\to$ ``long-term''
  \item ``limited shallower reefs'' $\to$ ``limited to shallower reefs''
  \item ``photodata'' $\to$ ``photo data''
  \item ``yielding to the complexity,'' ``proceeding section,''
    ``helps balances,'' ``use to model,'' ``the proposed pipelines ability,''
    ``proposed pipelines reproduces,'' ``may reflects'' --- all corrected
  \item ``Beijorn'' $\to$ ``Beijbom''
  \item ``Marcohorn'' $\to$ ``Macrohon''
  \item Informal acknowledgement rewritten: ``A massive thank you goes [to]
    Aidan Quigley'' $\to$ ``We thank Aidan Quigley for helping to initiate
    this collaboration.''
  \item SAM2 now correctly cites Ravi et al.\ (2024) (not the original SAM,
    Kirillov et al.\ 2023) at its first reference.
\end{itemize}
\end{response}

\bibliographystyle{apalike}
\bibliography{../references}

\end{document}
